%!TEX root = ../thesis.tex
% ******************************* Thesis Appendix B ********************************

\chapter{Interview Questions}\label{sec:interview}
The interview questions can be found below. The three interviewees are marked by their initials. Not all questions were answered by all interviewees.

ΕΠ: Δουλεύει σε έναν ξενώνα με χρόνιους ασθενείς εδώ και περίπου 20 χρόνια, είναι Διευθύντρια. 

ΚΝ: Δουλεύει ως νοσηλεύτρια σε Κοινοτική Ψυχιατρική Νοσηλευτική για περίπου 30 χρόνια.

ΣΤ: Δουλεύει για το Κέντρο Ψυχικής Υγείας, στην κινητή μονάδα Χανίων για 2,5 έτη. Πριν δούλευε σε ψυχογηριατρικό.


\textbf{Q1:	Ποιες είναι οι πιο συχνές ή χρονοβόρες εργασίες σε μια τυπική βάρδια;}

ΕΠ: Δύο είδη: Απλός νοσηλεύτης Vs. υπεύθυνος βάρδιας. 
Απλός νοσηλευτής: οι γυναίκες τις γυναίκες, οι άντρες τους άντρες.
Δεν χρεώνουν να έχεις τον τάδε ασθενή, ο κάθε νοσηλευτής πάει σε όποιον μπορεί.
Οι απλοί νοσηλευτές:
•	Δίνουν φάρμακα/αγωγή και ελέγχουν το στόμα ότι τα ήπιανε (κύριο μέλημα του νοσηλευτή!)
•	Δεν συνεργάζονται τόσο πολύ με τους ασθενείς για να τους διδάξουν τι έχουν – αυτό το κάνει ο γιατρός, γιατί δε λειτουργεί η ψυχοθεραπεία στον ψυχωσικό, ο νοσηλευτής δεν του πάει κόντρα, λέει «ό,τι είπε ο γιατρός θα κάνουμε, θα πιούμε αυτά τα χάπια», «θα περιμένουμε το γιατρό να μας πει πότε θα τα μειώσουμε/κόψουμε»
•	Παρακολουθεί τους ασθενείς: π.χ. αν έχει αϋπνίες δεν τους δίνει βραδινό φαγητό.
•	Να μη φύγει ο ασθενής, να μη κάνει απόπειρα αυτοκτονίας (αφαιρούν αναπτήρες, σουγιάδες, κλπ.)
•	Θα πρέπει να γνωρίζουν τους ασθενείς: αν είναι επικίνδυνοι, κλπ. Οι νέοι ασθενείς είναι πιο δύσκολοι. Συνδέονται και γίνονται και φίλοι, δηλ. τους παρηγορούν, ρωτάνε για την οικογένειά τους – ειδικά στους χρόνιους. Οι ασθενείς λένε πολλά πράγματα στους νοσηλευτές, ακόμα και που δεν λένε στους γιατρούς.
• Κανονίζει και τους γιατρούς για εξετάσεις, και τους συνοδεύουν: π.χ. αν θέλει οφθαλμίατρο, κλπ.
• Κάνουν καρδιογράφημα, αν δε νιώθει καλά.

ΚΝ: Είναι νοσηλεύτρια. Δουλεύει στην κινητή μονάδα του Κέντρο ψυχικής υγείας Χανιών - έχουν ιατρεία στο Νομό Χανίων με 380 ενεργά περιστατικά (κάνουν κυκλικά ιατρεία = follow up, έχουν θεραπευτική σχέση με τους ασθενείς), 5ήμερο - 8ωρο, μεγάααααλος φόρτος εργασίας, έχει θέση ευθύνης, συνέχεια καινούργια ΦΕΚ (που ζητάνε υποβολή εκθέσεων, στοχοθεσιών, αναφορών, φυλλάδια, έντυπα…). Ως τώρα τα έκανε μετά τις 3μμ, στο σπίτι της. Έχει ατζέντα γραπτή (όπου γράφει τα ενέσιμα, πορείες νόσου, κλπ.). Έχουν διασυνδέσεις με την κοινότητα για να μην απομονώνονται οι ασθενείς. Έχουν 11 περιοχές και κάθε περιοχή είναι χρεωμένη σε ένα κλιμάκιο, δηλαδή υπάρχουν σταθερά πρόσωπα αναφοράς.

ΣΤ: Εξυπηρετούν άτομα με νευρώσεις και ψυχώσεις, και κοινωνικά θέματα (γι’ αυτό έχουν και κοινωνικούς λειτουργούς ως συναδέλφους). Έχουν 380 ενεργά περιστατικά. Κάνουν 11 ιατρεία το μήνα, σε διάφορες περιοχές και χωριά (π.χ. στον Βάμο έχουν 70 ασθενείς).
Πώς δέχονται ασθενείς;
-	Τους παραπέμπουν ασθενείς οι γενικοί ιατροί ή το πρόγραμμα «Βοήθεια στο σπίτι».
-	Οι ίδιοι οι ασθενείς ή τα παιδιά τους.

Προβλήματα βάρδιας:
Δεν υπάρχουν ηλεκτρονικοί φάκελοι με πληροφορίες ασθενούς (π.χ. με τα εργαστηριακά του, τα παθολογικά/ψυχιατρικά του φάρμακα, κλπ.). Άρα οι νοσηλευτές πρέπει να συλλέγουν πληροφορίες από διάφορες πηγές, κι αυτό είναι χρονοβόρο.



\textbf{Q2:	Υπάρχουν ειδικές προκλήσεις που αφορούν μόνο την ψυχιατρική περίθαλψη κατά τη διάρκεια της νοσηλευτικής λογοδοσίας ή άλλων καθηκόντων;}
ΕΠ:
• Η ασφάλεια και να μη φύγουν μέχρι να δράσουν τα φάρμακα, όσο πιο γρήγορα δράσουν τόσο πιο γρήγορα θα συνέλθει.
• Να πάρει ο ασθενής τα φάρμακά του, να μην τα πετάξει.
• Να μην προκαλέσει κακό στον εαυτό του ή σε άλλον.

KN: Είναι πολύ σύνθετη δουλειά, γιατί ανοίγει μια μεγάλη «βεντάλια» με διαδικασίες.

ΣΤ: Ναι, ο σκοπός της δουλειάς είναι πολύπλευρος:
-	Συμμόρφωση με τα φάρμακα
-	Επεξήγηση στον ασθενή γιατί πρέπει να παίρνει τα φάρμακά του τακτικά (εξηγεί το θεραπευτικό πλάνο και στόχους)
-	Κοινωνικοποίηση των ασθενών και ένταξή τους στην κοινωνία
-	Παραμονή στην κοινότητα
-	Όχι υποτροπή ή μειωμένη υποτροπή

Η δουλειά της νοσηλεύτριας είναι πολύπλοκη: οι νοσηλευτές φροντίζουν τους ασθενείς όλο το μήνα, όχι μόνο στη διάρκεια των ιατρείων. Επίσης, έχουν και τακτική τηλεφωνική επικοινωνία για εκτίμηση της κατάστασής τους (κάνουν monitoring ανάλογα με τη διάγνωση, π.χ. ανάλογα με το τι έχει αλλάξει στον ασθενή μπορεί να αυξήσουν τη δόση κάποιου φαρμάκου. Ξέρουν ότι π.χ. στην νεύρωση γύρω στο μήνα θα πρέπει να έχει καλύτερη διάθεση ο ασθενής, κι αν δεν έχει, αυξάνουν τη δόση).


\textbf{Q3:	Με ποιον τρόπο παρακολουθούνται αυτή τη στιγμή οι εργασίες, οι σημειώσεις και οι θεραπείες των ασθενών;}
KN: ΌΛΑ υπήρχαν στο χαρτί, σε μπακαλοτέφτερα (μεγάλα τετράδια). Τώρα το κάθε τμήμα έχει το δικό του αρχείο. Θέλουν μια ενιαία βάση δεδομένων (για να έχουν στατιστικά, κλπ.). Εδώ και 1 χρόνο τούς έβαλε η Πληροφορική μέσα σε ένα πρόγραμμα, το ΑΜΕΤ (αμέτ): εκεί καταγράφονται όλα τα ραντεβού, τις συνταγογραφήσεις, τις αιμ/κές εξετάσεις των περιστατικών (είναι συνδεδεμένο με το νοσοκομείο).

Για κλείσιμο ραντεβού, πρέπει να υπάρχει αξιολόγηση από κοινωνικούς λειτουργούς, γιατί κάποιο περιστατικό μπορεί να είναι επείγον (π.χ. πρόθεση για αυτοκτονία ή ανθρωποκτονία)!

ΣΤ: Γραπτά σε φακέλους. Ηλεκτρονικά μόνο τι φάρμακα έχουν συνταγογραφηθεί και έχει πάρει ο ασθενής.

\textbf{Q4:	Πώς ιεραρχείτε και παρακολουθείτε τις εργασίες κατά τη διάρκεια της βάρδιας σας;}

\textbf{Q5:	Πόσο συχνά χρειάζεται να ενημερώνετε ή να ανακτάτε πληροφορίες για τον ασθενή; Και ποια είναι η διαδικασία;} 

ΕΠ: Διαφορετικά τα οξέα με τα χρόνια, γιατί στα χρόνια τον γνωρίζουν τον ασθενή.

ΣΤ: Σε κάθε τηλεφωνική επικοινωνία ή ιατρείο χρειάζεται να βλέπει κάποιες πληροφορίες (π.χ. τα αποτελέσματα των εργαστηριακών εξετάσεων) για να κάνει monitor της εξέλιξης της υγείας/θεραπείας. Έχει ένα excel sheet όπου γράφει τις προσωπικές της σημειώσεις για κάθε ασθενή, για να μη χρειάζεται να ανατρέχει και να ξεφυλλίζει τον μοναδικό, γραπτό και πολυσέλιδο Φάκελο Ασθενούς κάθε φορά… Π.χ. γράφει το ΑΜΚΑ (για να δει στο σύστημα η ίδια αν έχει ή δεν έχει φάρμακα) + τηλέφωνο ασθενή, για την αγωγή του ασθενούς, πότε χρειάζεται ο κάθε ασθενής την επόμενη συνταγογράφηση, πότε έκανε ενέσιμο και πότε πρέπει να κάνει το επόμενο, πότε έκανε αιματολογικές και πότε πρέπει να κάνει τις επόμενες. Δε χρειάζεται να γράφει τα συμπτώματα, γτ της είναι γνωστά. Η διάγνωση είναι δυναμική, δηλαδή μπορεί να αλλάξει (π.χ. ένας ασθενής μπορεί να είναι γραμμένος με κατάθλιψη, αλλά τελικά να έχει μανιοκατάθλιψη, γτ του δίνεις χάπια κατάθλιψης και «πετάει»!). Επίσης, γράφει κάποιο SOS σχόλιο από τους γιατρούς, π.χ. ότι το τάδε φάρμακο μπορεί να φέρει αϋπνίες. Γράφει και για τα επιδόματα που έχει πάρει ο ασθενής (π.χ. ΚΕΠΑ), κι αν έχει ποσοστό αναπηρίας (αυτά ενδιαφέρουν τον κοινοτικό νοσηλευτή, όπως η ίδια), γτ οι περισσότεροι ασθενείς δεν μπορούν να εργαστούν και ίσως δεν ξέρουν ότι δικαιούνται επιδόματα για να ζήσουν. Την ενδιαφέρει επίσης αν ο ασθενής έχει ασφάλιση, αν έχει τα προς το ζην, αν έχει οκ ή προβληματική οικογένεια, αν η οικογένειά του αποδέχεται τον άρρωστο (αν τον αποδέχεται, τότε θα κάνει λιγότερες υποτροπές ή αν η μητέρα του έχει υπερεμπλοκή/αν έχει τραύμα από κακοποίηση, τότε θα έχει περισσότερες υποτροπές).

\textbf{Q6:	Υπάρχει κάποια βάση δεδομένων; Ποιο σύστημα καταγραφής χρησιμοποιείται;}

ΕΠ: Παλιά: Φάκελος Ασθενούς χειρόγραφα, με δημογραφικά, κοινωνικά (τι δουλειά κάνει, έχει παιδιά, αδέρφια;), ιατρικό ιστορικό τους, έχει κοινωνικό δίκτυο, τηλέφωνο ανάγκης (ποιος έχει την επιμέλειά του), πόσες φορές έχει μπει, έχει βγει και πότε, ποιος γιατρός τον είχε δει.
Η κοινωνική λειτουργός παίρνει το Κοινωνικό Ιστορικό, το οποίο συμπληρώνει μετά ο γιατρός με το ιατρικό ιστορικό (ακούει φωνές, έχει ψευδαισθήσεις, τι συμπτώματα έχει δηλαδή). Μετά υπάρχει και επισκέπτρια υγείας που καταγράφει πόσες φορές πηγαίνει στο νοσοκομείο, με ποιον έχει συνομιλίες, πόσες φορές έχει πάει στο σπίτι του, αν έχει νοσηλευτεί ο ασθενής, όταν ξανάρθει ο ασθενής αν έχει προκύψει ένα χειρουργείο ή μια ασθένεια. Και η ψυχολόγος, τι του αρέσει, τι τον φοβίζει, τι θα ήθελε, με ποιον του αρέσει να συζητάει και ποιοι είναι οι φίλοι του, ποιο είναι ένα πρόσωπο αναφοράς, και τεστ προσωπικότητας, και βγαίνει ένα προφίλ, δηλαδή ποιες είναι οι ανάγκες του.
Μετά όλοι αυτοί μοιράζονται τις πληροφορίες προφορικά, δηλαδή μαζεύονται καθημερινά σε μια αίθουσα και κάνουν συνάντηση προσωπικού. Στα ΟΞΕΑ γίνεται αυτό. Στα ΧΡΟΝΙΑ γίνονται πιο αραιά οι συναντήσεις. Τα καταγράφουν γραπτά, με το χέρι, στον Φάκελο Ασθενούς [είναι 1 μόνο, δηλαδή αν χαθεί, είναι πρόβλημα] οι κοινωνικοί λειτουργοί.
Το μόνο ηλεκτρονικό σύστημα: πλέον με το ΑΜΚΑ φαίνονται κάποια πράγματα – φαίνεται η αγωγή του και σε ποιον γιατρό πήγε, και πόσες φορές μπήκε ή βγήκε.
Άλλο που υπάρχει – γραπτά – είναι ιατρική γνωμάτευση (αν μπορεί να εργαστεί, κλπ.) στο αρχείο της Γραμματείας. Έχει συγκεκριμένη φόρμα.

KN: Μόνο για 1 χρόνο, δηλαδή τώρα τελευταία. Έχουν πολύ μεγάλη ανάγκη από την τεχνολογία, έχουν κουραστεί να γράφουν βιβλία, π.χ. ποιοι προσήλθαν στα ιατρεία – αν και τώρα αυτό αντικαθίσταται από το ΑΜΕΤ.


\textbf{Q7:	Πώς αποθηκεύονται τα δεδομένα;} 

EΠ: Παλιά χειρογραφα, τώρα ηλεκτρονικά.
ΚΝ: Παλιά όλα χειρόγραφα, τώρα μερικά ηλεκτρονικά.

\textbf{Q8:	Τι θα θέλατε να βελτιώσετε από τον υφιστάμενο τρόπο παρακολούθησης;}
ΣΤ: Θα βοηθούσαν πολύ τα διαγράμματα (π.χ. πόση κατάθλιψη έχει ο ασθενής από το 0 ως το 10), οι φόρμουλες, τα templates (π.χ. Τι φωνές ακούς; Α. σαν από ραδιόφωνο; Β. μιλάνε η μία στην άλλη; κλπ.) για να μπορούν να εξάγονται μετρήσεις και στατιστικά για τους ασθενείς.

\textbf{Q9:	Ποιο είδος πληροφοριών είναι απαραίτητο κατά τη διάρκεια της λογοδοσίας; Μπορείτε να αναφέρετε τις βασικότερες κατηγορίες πληροφοριών;}

ΕΠ: Τα γράφουν κάτω, τι κάνανε και τι πρέπει να γίνει (τι υλικά, χάπια χρειάζονται) για την επόμενη βάρδια. Συνήθως τη γράφει ένας: ο υπεύθυνος βάρδιας ή ο προϊστάμενος.
Υπάρχει και Φάκελος Ασθενούς στα πανεπιστημιακά νοσοκομεία.
Ενημερώνεις σε ποιον άλλο γιατρό πήγε ένας ασθενής, αν έφαγε, δεν έφαγε, αν πήρε τα φάρμακά του ή όχι, και αφήνεις εκκρεμότητες.

ΚΝ: Πού μπορεί να βοηθήσει το ΑΙ;
-	για την καταγραφή της ημέρας (προς το παρόν τα γράφει πολύ συντομευμένα για να βγάλει στατιστικά, γτ το Υπουργείο τής ζητάει στατιστικά), π.χ. ότι έκανε τηλεφωνική επικοινωνία με τον τάδε, κλπ. (μέχρι τώρα όλα αυτά σε ένα τεφτέρι!) για κατηγοριοποίηση των τηλεφώνων και λοιπών υπηρεσιών.
Ποιες είναι οι υπηρεσίες;
Α) πορεία νόσου
Β) τηλεφωνική επικοινωνία για την φαρμακευτική του αγωγή 
Γ) διασύνδεση με φορείς (π.χ. με Μητρόπολη, με σχολεία [όλων των επιπέδων] για Πρόγραμμα Αγωγής Υγείας, με παιδικούς σταθμούς για Πρώιμη Ανίχνευση, κέντρα κοινότητας, π.χ. για προγράμματα βοήθειας στο σπίτι) και υπηρεσίες (π.χ. με Δήμους, με της Αστυνομίας το Τμήμα Κακοποίησης, με την Εισαγγελία…), γτ η δουλειά τους βασίζεται στο δίκτυο συνεργασίας με όλες αυτές τις υπηρεσίες. Καταρχάς, τα ιατρεία της τα κάνει σε Κέντρα Υγείας ή σε περιφερειακά ιατρεία ως φιλοξενούμενη ή τα προγράμματα Δήμου μπορεί να τα κάνει ακόμα και σε χώρο της Εκκλησίας. ΟΛΑ ΑΥΤΑ ΚΑΤΑΓΡΑΦΟΝΤΑΙ ΣΕ ΤΕΤΡΤΑΔΙΑ.
Το πρώτο βήμα: Τώρα αυτά είναι σε ΕΝΑ ΜΟΝΑΔΙΚΟ τετράδιο με 4 στήλες, το λένε «ΗΜΕΡΗΣΙΑ ΚΙΝΗΣΗ» (αντί για ΛΟΓΟΔΟΣΙΑ ΠΟΥ ΕΙΝΑΙ ΟΥΣΙΑΣΤΙΚΑ), και γράφουν με απλές λέξεις-κλειδιά για να κερδίσουν χρόνο.

Το δεύτερο βήμα, μετά την ΗΜΕΡΗΣΙΑ ΚΙΝΗΣΗ:
ΦΑΚΕΛΟΣ ΑΣΘΕΝΟΥΣ: ένας για κάθε ασθενή, χειρόγραφος.
ΓΡΑΦΟΥΝ αναλυτικά τι είχε η τηλεφωνική επικοινωνία, ή αν τα φάρμακα είχαν παρενέργειες, αν κοιμάται, αν τρώει, τι της περιέγραψε, τις της είπε, τι διάθεση είχε… -> όλα αυτά έχουν σημασία για το γιατρό!
-	Εργαστηριακές εξετάσεις,
-	Κοινωνικό ιστορικό (αμκα, δημογραφικά στοιχεία, υποστηρικτικό πλαίσιο, κλπ. )
-	Ιατρικό ιστορικό (η νόσος)
-	Παρατηρήσεις (π.χ. υπάρχει μια διέγερση, διαταραχή στη σκέψη)
Υπάρχει και το Προσωπικό Αίτημα (Intake): η πρώτη επαφή με το περιστατικό είναι πολύ σημαντική, προχωράς βήμα βήμα και μπορεί να βοηθήσουν τα νοσηλευτικά πρωτόκολλα!!!
«Θέλω βοήθεια, δεν είμαι καλά» -> τηλεφωνικά ο νοσηλευτής εκμαιεύει
Έχει «Έντυπο Καταγραφής Διαδικασιών» -> δείχνει τι κάνουν και πού τα καταγράφουν ήδη!

\textbf{Q10:	Χρησιμοποιούνται συγκεκριμένα νοσηλευτικά πρωτόκολλα;}

ΕΠ: Βιολογικά (είναι καθαρός;),
Η σκέψη (είναι καθαρή;) ο λόγος (κατανοητός;).
Έλεγχος για επικίνδυνα αντικείμενα.
Οξεά τμήματα: Τα ενδιάμεσα είναι τα Κέντρα Ψυχικής Υγείας, όπου γίνεται μια νοσηλεία διάρκειας περίπου 15-20 ημερών. Μετά οι ασθενείς επιστρέφουν στο σπίτι τους, με ή χωρίς οικογένεια, και μόνο όταν τους ξαναπιάσει κρίση επιστρέφουν στην κλινική ή επισκέπτονται περιοδικά την ψυχιατρική κλινική για τη συνταγογράφηση των φαρμάκων.
Ο γιατρός έρχεται και παίρνει το ιστορικό, ενώ ο νοσηλευτής δεν ρωτάει πολλά. Σήμερα, τα οξεά περιστατικά αντιμετωπίζονται αποκλειστικά στα οξεά τμήματα της ψυχιατρικής κλινικής.

Χρόνια τμήματα: Αφορούν μόνιμους ασθενείς που διαμένουν σε ξενώνες ή διαμερίσματα, υπό την ευθύνη του Κέντρου Ψυχικής Υγείας. Οι νοσηλευτές τούς γνωρίζουν καλά και δεν χρειάζεται να παίρνουν κάθε φορά νέο ιστορικό. Στους χρόνιους ασθενείς, στους ξενώνες, δεν χρειάζεται κάποιος βοηθός για να υπενθυμίζει το ιστορικό, γιατί είναι μόνο δώδεκα ασθενείς όλοι κι όλοι, και τους γνωρίζουν ήδη.

Ο φάκελος θα μπορούσε να υπενθυμίζει, για παράδειγμα, πώς γίνεται η αλλαγή φλεβοκαθετήρα ή ουροκαθετήρα, ή πώς λαμβάνεται αίμα — κάτι ιδιαίτερα χρήσιμο για νέους νοσηλευτές.
Το τελετουργικό, δηλαδή η ίδια η διαδικασία, δεν πραγματοποιείται συχνά, και οι νοσηλευτικές ενέργειες μερικές φορές ξεχνιούνται.


\textbf{Q11:	Υπάρχει κάποια προηγούμενη εμπειρία με την τεχνολογία στη ροή εργασίας σας, την οποία θεωρήσατε ιδιαίτερα χρήσιμη ή μη χρήσιμη; /	(Πώς) έχετε προσπαθήσει να λύσετε το πρόβλημα της χειρόγραφης καταγραφής φακέλου ασθενούς; / Τι κάνετε ως τώρα για να βελτιώσετε την κατάσταση; } 

ΕΠ: Παλαιότερα δεν υπήρχε ενδοεπικοινωνία, τίποτα· συνεννοούνταν μόνο με τηλέφωνα.
Πλέον χρησιμοποιούν κινητά, άλλωστε οι κλινικές είναι μικρές.
Υπήρχαν βομβητές στα πανεπιστημιακά νοσοκομεία, αλλά αυτό γινόταν παλιά.

Ο νοσηλευτής είναι συνεχώς απίκο στις ψυχιατρικές κλινικές, ώστε να παρακολουθεί τι κάνουν οι ασθενείς.
Οι συγγενείς έρχονται μόνο για επίσκεψη — δεν κοιμούνται ούτε παραμένουν μαζί με τους ασθενείς.

ΚΝ: Έχει χρησιμοποιήσει ΑΙ (chatgpt μάλλον) για τις αναφορές και τις εκθέσεις και περιλήψεις, π.χ. για την Ευρωπαϊκή Επιτροπή. => για γραφειοκρατία το θέλει.
Ήταν επιστημονική φαντασία ο υπολογιστής όταν ήταν νέα, είναι 56 χρονών τώρα! Η Διευθύντριά της είναι 70 χρονών. Αποφεύγουν την πολλή χρήση με τον υπολογιστή. Αλλά έχει νέα παιδιά στην δουλειά που είναι πρόθυμα να χρησιμοποιήσουν ΑΙ.
Τι της αρέσει: που δεν είναι ξένο προς αυτόν που το ζητάει (το αποτέλεσμα του ΑΙ).
Μπορεί να βοηθήσει στα γραφειοκρατικά, στα στατιστικά κλπ., αλλά όχι στη νοσηλεία ή στην επαφή με τον ασθενή.
μπορεί μόνο να δώσει έναν μπούσουλα, π.χ. «πότε μπορώ να κάνω γενική αίματος?»
Μια συνάδελφός της έχει γράψει πληροφορίες για τους ασθενείς της στο κινητό της, αλλά δεν ξέρει τι ακριβώς.

ΣΤ: Έχει ένα excel sheet όπου γράφει προσωπικές της σημειώσεις για κάθε ασθενή, για να μη χρειάζεται να ανατρέχει και να ξεφυλλίζει τον μοναδικό, γραπτό Φάκελο Ασθενούς κάθε φορά. 

\textbf{Q12:	Ποια θα ήταν τα πιο πολύτιμα χαρακτηριστικά σε έναν φωνητικό βοηθό για την εργασία σας;}

ΣΤ: Αμεσότητα, δηλαδή το να βρίσκει αμέσως την πληροφορία που ψάχνει χωρίς να πρέπει να ανατρέξει σε διάφορες άλλες πηγές.


\textbf{Q13:	Αισθάνεστε άνετα με τη χρήση φωνητικών εντολών σε ένα σύστημα; Αν όχι, ποιοι είναι οι λόγοι που σας αποτρέπουν;}
ΣΤ: Δεν είπε ξεκάθαρα ότι αισθάνεται άνετα, αλλά είπε ότι θα της άρεσε και θα το ήθελε να μπορεί να ρωτάει με φωνή ερωτήσεις για τους ασθενείς, π.χ. «Τι αγωγή παίρνει στην παρούσα φάση;»

\textbf{Q14:	Θα προτιμούσατε ο βοηθός να έχει έναν πιο ενεργητικό ρόλο (να σας κάνει προληπτικά υπενθυμίσεις, να αναφέρει στοιχεία ακόμα και αν δεν απευθύνεστε σε αυτόν, να είναι πιο συνομιλιακός) ή να λειτουργεί επικουρικά (να ξεκινάτε εσείς αποκλειστικά την αλληλεπίδραση);}
ΕΠ:
•	Φωνητική καταγραφή για ό,τι θέλει να πει, ώστε να μην τα ξεχάσει, επειδή συμβαίνουν διάφορα πράγματα και δεν προλαβαίνει να τα σημειώσει εκείνη τη στιγμή — για να τα καταγράψει αργότερα.
• Να ξέρει την αγωγή του.
• Να ξέρει πότε ήρθε τελευταία φορά.
• Να γνωρίζει λίγο τη γενική του εικόνα.
• Αν είχε παρενέργειες από κάποιο φάρμακο (ποιο φάρμακο του προκάλεσε τι) και γιατί το άλλαξε.
•	Σαν ένα ΠΡΟΧΕΙΡΟ ΠΡΟΣΩΠΙΚΟ μικρό φακελάκι ασθενούς.
Με ποιο μέσο: Κατά κύριο λόγο να είναι γραμμένα (γτ είναι οπτικός τύπος κυρίως), αλλά και να τα ακούει κιολας, αν θέλει. Όχι μόνο ακουστικά γιατί τα ξεχνάει. Για μια σύντομη πληροφορία μόνο, τότε οκ η φωνή.

ΚΝ: Patient Data – είναι εξατομικευμένο, δεν μπορείς να το εφαρμόσεις παντού, αλλά υπάρχει μια φόρμα με 5 βασικά πράγματα:
-	Επίπεδα φαρμάκων, που είναι μετρήσιμα, θα τους βοηθάει να μαθαίνουν αμέσως τα επίπεδα
-	Έχει βλεμματική επαφή;
-	Είναι φροντισμένος;
-	Είναι διεγερτικός;
-	Κάνει γλωσσική σαλάτα;
-	Έχει παραλήρημα;
Και μεταφέρει τα συμπτώματα στο γιατρό.

Ασχολούνται με τις παρακάτω περιπτώσεις ασθένειας:
-	Νεύρωση: θέλει όρια
-	Ψύχωση / Μανιοκατάθλιψη: θέλει ντάντεμα
1.	Παίρνει φάρμακα;
2.	Αν ναι, χτίζω δεξιότητες, π.χ. αυτοφροντίδας (δεν πλένονται)
3.	Φαρμακευτική αγωγή (παρενέργειες: τάση για έμετο, υπνηλία, μειωμένη πίεση…), Αυτοφροντίδα, Κοινωνική Απόσυρση, Δεξιότητες Πνευματικές-Σωματικές, Γενική Αίματος
-	Άνοια:
-	Διαχείριση Πένθους:

ΣΤ: Θα τη βοηθούσε να είχε υπενθυμίσεις και στατιστικά (π.χ. πότε υποτροπιάζει η διπολική διαταραχή).

\textbf{Q15:	Πώς αισθάνεστε σχετικά με τη χρήση ενός φωνητικού βοηθού όσον αφορά το απόρρητο των ασθενών και την ασφάλεια των δεδομένων; Υπάρχουν συγκεκριμένες ανησυχίες σχετικά με τις αλληλεπιδράσεις των ασθενών που θα πρέπει να ληφθούν υπόψη για τις φωνητικές εγγραφές ή τις απομαγνητοφωνήσεις;}
ΕΠ: Ναι, πού θα φυλάσσεται αυτό το υλικό, για να μην το κλέψει κάποιος από το κινητό του νοσηλευτή; Είναι ένα μεγάλο θέμα τα Ευαίσθητα Δεδομένα.


\textbf{Q16:	Πώς θα θέλατε να δίνετε ανατροφοδότηση σχετικά με τις αποκρίσεις του βοηθού; Kέιμενο ή φωνή;} 
ΕΠ: Και τα δυο!


----------------------------------------------------------------------

Επιπρόσθετες πληροφορίες για τις ελληνικές ψυχιατρικές δομές:

Ψυχιατρεία (εξαρτημένοι και ελεγχόμενοι ασθενείς· τρώνε μόνο με κουτάλια και παίρνουν τα χάπια τους. μπαίνουν εκεί οξέα περιστατικά και μένουν επίσης όσοι η οικογένειά τους δεν θέλει να τους στηρίξει. έχει πολλά τμήματα, περίπου 30 ασθενείς το καθένα, συνολικά γύρω στα 300 άτομα).

Εξωτερικά ιατρεία / κέντρα ψυχικής υγείας ή κινητές μονάδες (ο νοσηλευτής πηγαίνει κατ’ οίκον ή εργάζεται στα «νοσοκομεία ημέρας» – μια μεγάλη αίθουσα, όχι ακριβώς νοσοκομείο – όπου καλούνται οι ασθενείς κάθε πρωί να πηγαίνουν, να κάνουν δραστηριότητες και ο νοσηλευτής τούς καθοδηγεί και τους δίνει φάρμακα. το μεσημέρι φεύγουν και πάνε σπίτι τους).

Ψυχιατρικές κλινικές των πανεπιστημιακών νοσοκομείων (σε κάθε νοσοκομείο υπάρχει τμήμα οξέων περιστατικών· ο νοσηλευτής είναι συνέχεια μαζί με τον ασθενή, από τη στιγμή που θα έρθει, μέχρι – αν χρειαστεί – να απομονωθεί ή να δεθεί, έως ότου δράσει το φάρμακο. υπάρχει συστηματική παρακολούθηση).

Ψυχιατρικές κλινικές γενικά: τα πράγματα είναι λίγο πιο ανθρώπινα – οι ασθενείς κοιμούνται εκεί, τρώνε κλπ.

Ξενώνες: λειτουργούν μετά τα ψυχιατρεία, για να αυτονομηθούν οι ασθενείς. φιλοξενούν 12-15 άτομα και έχουν μια σταθερή ρουτίνα: πρωινό στην κουζίνα ή τραπεζαρία (με πιρούνια κανονικά), μπάνιο κ.ά.

Διαμερίσματα: εκεί μένουν όπως οι φυσιολογικοί άνθρωποι.

Στο ψυχιατρείο παλαιότερα έμπαιναν και άτομα με νοητική υστέρηση, δηλαδή με άλλα ψυχοπαθολογικά συμπτώματα.

Ψύχωση: οι ασθενείς ακούνε φωνές ή έχουν ακουστικές ψευδαισθήσεις και συχνά σταματούν να παίρνουν τα φάρμακά τους (επειδή νομίζουν ότι τα φάρμακα τούς τρελαίνουν, ενώ στην πραγματικότητα τους βοηθούν, ή επειδή έχουν παρενέργειες). Τότε οι φωνές επιστρέφουν και γι’ αυτό μπαινοβγαίνουν στα νοσοκομεία. Το ΠΡΩΤΟ πράγμα που πρπ να  καταλάβει ο ασθενής είναι να ΠΑΙΡΝΕΙ ΤΑ ΦΑΡΜΑΚΑ, γτ χάνει δεξιότητες.
Τα φάρμακα όμως δεν αντιμετωπίζουν την έλλειψη βούλησης ή την κοινωνική απομόνωση.








